% vim: spl=pt
\begin{exercício}{Espectro de um operador unitário}{ex14}
    Mostre que, se \(U\) é um operador unitário, então
    \begin{enumerate}[label=(\alph*)]
        \item \(U\) é uma isometria sobrejetiva (bijetiva); e
        \item \(\sigma(U) \subset \setc{z \in \mathbb{C}}{\abs{z} = 1}\).
    \end{enumerate}
\end{exercício}
\begin{proof}[Resolução de (a)]
  Como \(U\) é unitário, temos \(U \in \bounded(\hilbert)\) e \(U^* = U^{-1}.\) Assim, \(U\) é bijetivo e, em particular, sobrejetivo. Notemos que \(U\) preserva o produto interno, já que
  \begin{equation*}
    \forall x,y \in \hilbert: \inner{Ux}{Uy} = \inner{x}{U^*Uy} = \inner{x}{y}.
  \end{equation*}
  Assim, temos
  \begin{equation*}
    \forall x \in \hilbert: \norm{Ux} = \sqrt{\inner{Ux}{Ux}} = \sqrt{\inner{x}{x}} = \norm{x},
  \end{equation*}
  logo \(U\) é uma isometria.
\end{proof}
\begin{proof}[Resolução de (b)]
  Como \(U\) é unitário, é uma isometria bijetiva, portanto \(\norm{U} = 1\). Ainda, \(U\) é normal já que \(U^* U = U U^*,\) portanto \(r(U) = \norm{U} = 1\) e concluímos que \(\sigma(U) \subset D,\) onde \(D \subset \mathbb{C}\) é o disco fechado unitário. Temos o mesmo resultado para \(U^* = U^{-1}\), isto é, \(\sigma(U^{-1}) \subset D\). Disso segue que
  \begin{align*}
    \lambda \in \sigma(U) &\iff \lambda \in \sigma(U) \land \lambda^{-1} \in \sigma(U^{-1})\\
    &\implies\abs{\lambda} \leq 1 \land \abs{\lambda}^{-1} \leq 1\\
    &\iff \abs{\lambda} \leq 1 \land \abs{\lambda} \geq 1\\
    &\iff \abs{\lambda} = 1,
  \end{align*}
  portanto \(\sigma(U) \subset \setc{z \in \mathbb{C}}{\abs{z} = 1}.\)
\end{proof}
